\documentclass[10pt,a4paper]{amsart}
\usepackage[margin=19mm]{geometry}
\usepackage{amsmath,amssymb,mathtools}
\usepackage[scr=rsfs,cal=euler]{mathalfa}
\usepackage{xfrac}
\usepackage[unicode,hidelinks,bookmarks=false]{hyperref}
\newcommand{\ie}{i.e.\ }
\newcommand{\cf}{cf.\ }
\newcommand{\resp}{respectively\ }
\newcommand{\Subsection}{\subsection}
\providecommand{\nobreakdash}{\nobreak\hbox{-}}
\setlength{\emergencystretch}{2em}
\multlinegap0pt
\usepackage{amssymb}
\usepackage{bbm}
\usepackage{enumerate}
\usepackage[all]{xy}
\usepackage{float}
\usepackage{cancel}
\usepackage{csquotes}
\newtheorem{lemma}{Lemma}
\newtheorem{theorem}{Theorem}
\DeclareMathOperator{\Aut}{Aut}
\newcommand{\G}{\Gamma}
\newcommand{\Ra}{\Rightarrow}
\newcommand{\lra}{\longrightarrow}
\newcommand{\oAut}{o\text{-}\Aut}
\newcommand{\ra}{\rightarrow}
\newcommand{\vAut}{v\text{-}\Aut}
\title{Automorphisms of valued Hahn groups}
\author{Salma Kuhlmann, Michele Serra}
\date{Source: arXiv:2302.06290v4 (CC BY 4.0)}
\begin{document}
\maketitle

\noindent\emph{Source and licence.} Automorphisms of valued Hahn groups. Version arXiv:2302.06290v4 (CC BY 4.0), distributed by arXiv under the Creative Commons Attribution 4.0 International licence, http://creativecommons.org/licenses/by/4.0/, as recorded in the arXiv licence metadata for this exact version and shown on the abstract page https://arxiv.org/abs/2302.06290v4. Full licence notice: \texttt{licenses/arxiv-2302.06290-CC-BY-4.0.txt}.\par\smallskip

\noindent\textit{Editorial scope.} Counted unit: the article's theorem lifting-principal-convex-sbgps (source0.tex lines 593--605). The article's complete original proof of it is delivered with it (source0.tex lines 606--647, one \texttt{proof} environment of 42 lines). No mathematics is written, repaired or re-derived: the statement and the proof are the article's own lines, in the article's own notation, and no retained line is substituted or re-flowed.  Retained uncounted as the source-local context the proof needs: lines 572--591.  Nothing was omitted from any retained span, so the package carries no omission record.  No retained line contains a citation, so the wrapper carries no bibliography and no bibliography entry.  No retained line contains a figure environment, an image-inclusion command or drawing code, so no image is delivered and the package declares no asset dependency.  Editorial formatting only, in addition: an AMS \texttt{amsart} preamble that restates the article's own declarations and macros used by the retained lines, namely the environments \texttt{lemma}, \texttt{theorem} on separate counters, together with the standard packages those lines need.  The retained environments print in this document as Lemma 1, Theorem 1 in order of appearance, whereas the article prints the same items with the numbering of its own full document.  The complete native source of the article as distributed in the arXiv e-print for this exact version is bundled unchanged as \texttt{sources/137313/source0.tex}.
\begin{lemma}
	\label{h-commutes-C}
	Let $\sigma\in\vAut G$ and let $g\in G$ and $C_g := \{x\in G:\ v(x)\geq v(g)\}$. Then we have:
	$$\sigma(C_g) = C_{\sigma(g)} = \{x:\ v(x) \geq v(\sigma(g))\}.$$
	\begin{proof}
		Let $ \sigma_\Gamma$ be the automorphism of the chain $\Gamma$ induced by $\sigma$.
		Let $x\in C_g$. Then
		$$v(x)\geq v(g) \Ra v(\sigma(x)) = \sigma_\Gamma (v(x)) \geq \sigma_\Gamma (v(g)) = v(\sigma(g))$$
		so $\sigma(x) \in C_{\sigma(g)}$, that is $\sigma(C_g) \subseteq C_{\sigma(g)}$.
		
		Vice versa, let $x\in C_{\sigma(g)}$. Then
		$$
		v(x) \geq v(\sigma(g)) \Ra
		v(\sigma^{-1}(x)) = \tilde \sigma^{-1}(v(x)) \geq \sigma_\Gamma^{-1}(v(\sigma(g)))= v(\sigma^{-1}(\sigma(g)))=v(g)
		$$
		because $\sigma_\Gamma^{-1}$ is also an automorphism of $\Gamma$.
		So $\sigma^{-1}(x) \in C_g$, that is $x\in \sigma(C_g)$.
		Hence $\sigma(C_g) \supseteq C_{\sigma(g)}$, which completes the proof.
	\end{proof}
\end{lemma}

\begin{theorem}
	\label{lifting-principal-convex-sbgps}
	Let $\tau\in\Aut(\G, <)$. Then $\tau$ lifts to an automorphism
	$\sigma \in \oAut G$ if and only if, for all $\gamma\in\Gamma$ there is an isomorphism
	$$
	\sigma_\gamma\colon C_\gamma \overset{\sim}{\lra} C_{\tau(\gamma)}
	$$
	such that 
	\begin{equation}
	\label{cond-gamma-delta}
	\gamma>\delta \Rightarrow \sigma_\delta|_{C_\gamma}=\sigma_\gamma.
	\end{equation}
\end{theorem}

\begin{proof}
	Assume that $\tau$ lifts to $ \sigma\in\oAut G $. Let $\gamma\in\Gamma$ and choose $g\in G$ such that
	$v(g) = \gamma$. Recall that, by definition, $v(\sigma(g)) = \tau (v(g))$.
	Then Lemma~\ref{h-commutes-C} gives
	$$
	C_\gamma =C_g = \{x:\ v(x)\geq \gamma = v(g)\}\simeq C_{\sigma(g)} = \{x:\ v(x) \geq \tilde \sigma(v(g)) = \tau(\gamma)\} = C_{\tau(\gamma)}.
	$$
	where the isomorphism is simply given by $\sigma$, which is order preserving, so
	condition \eqref{cond-gamma-delta} is satisfied.
	
	Vice versa, suppose that for all $\gamma\in\Gamma$ there is an isomorphism
	$$
	\sigma_\gamma\colon C_\gamma \overset{\sim}{\lra} C_{\tau(\gamma)}
	$$
	such that $\gamma>\delta \Rightarrow \sigma_\gamma(C_\gamma)\subset \sigma_\delta(C_\delta)$.
	We want to construct a lift $\sigma \in \oAut G$ of $\tau$. We define it as follows:
	for $x\in G$ with $v(x) = \gamma$, set $\sigma(x):=\sigma_\gamma(x)$. First let us show that
	$\sigma$ is a well defined group morphism. Let $x,y \in G$. We have two cases:\\
	\textbf{Case 1.}
	$v(x)\neq v(y)$, and we can assume $v(x)<v(y)$.
	Then $v(x+y)=v(x) =:\gamma$. So we have 
	$$
	\sigma(x-y) = \sigma_\gamma (x-y) = \sigma_\gamma(x) - \sigma_\gamma(y) = \sigma(x) - \sigma_{v(y)}(y) = \sigma(x) - \sigma(y)$$
	where in the second-last equality we used condition~\eqref{cond-gamma-delta}.\\
	\textbf{Case 2.} $v(x) = v(y) = \delta$. Then $v(x+y) =: \gamma \geq \delta$. Then
	$$
	\sigma(x-y) = \sigma_\gamma(x-y) = \sigma_\delta(x-y) = \sigma_\delta(x) - \sigma_\delta(y) = \sigma(x) - \sigma(y).
	$$
	Again, for the second equality we used condition~\eqref{cond-gamma-delta}.
	
	To show that $\sigma$ is an automorphism, we define its inverse. Similarly to what we did
	to define $\sigma$ we define $\sigma^{-1}$ as follows: for $x\in G$ with 
	$v(x) = \tau(\gamma)$ we set $\sigma^{-1}(x) = \sigma_\gamma^{-1}(x) $.
	Checking that this is a
	well defined group morphism is done the same way as it was done for $\sigma$. To see that this is indeed the inverse of $ \sigma $, let $x\in G$ with 
	$v(x)=\gamma$.
	Recall that $\sigma_\gamma(x) \in C_{\tau(\gamma)}$ and $\sigma_\gamma^{-1}: C_{\tau(\gamma)}\ra C_\gamma$. So
	$$
	\sigma^{-1}(h(x)) = \sigma^{-1}(\sigma_\gamma(x)) = \sigma_\gamma^{-1}(\sigma_\gamma(x)) = x
	$$
	and similarly $\sigma(\sigma^{-1}(x)) = x$. So $\sigma$ is an automorphism of $G$, which completes the proof.
\end{proof}

\end{document}
